\documentclass[12pt]{article}

\usepackage[american]{babel}
\usepackage[margin=1in]{geometry} 
\usepackage{amsmath,amsthm,amssymb}
\usepackage{multicol}
\setlength{\columnsep}{1cm}
\usepackage{mathrsfs}
\usepackage{tikz}
\usepackage{color}
\usepackage{lipsum} % for filler text
\usepackage{fancyhdr}
\usepackage{graphicx}
\pagestyle{fancy}
\usepackage[svgnames]{xcolor}
\usepackage[colorlinks=true, allcolors=DeepPink]{hyperref}
\usepackage{hyperref}
\usepackage{float}
\usepackage{subcaption}
\usepackage{xcolor}
\usepackage{wrapfig}

\newcommand{\N}{\mathbb{N}}
\newcommand{\Z}{\mathbb{Z}}
\newcommand{\ep}{\varepsilon}
\newcommand{\pa}{\partial}

\usepackage{titling}
\usepackage{hyperref}
\usepackage{longtable}

\setlength{\droptitle}{-8em}   % This is your set screw
\setlength{\headheight}{16pt}

\begin{document}

\title{AMATH 732 Project Ideas}
\author{the royal you (you as in we)}
\maketitle

\section{Introduction}
i love latex
yay $\epsilon$ $\varepsilon$

\begin{table}[H]
    \centering
    \begin{tabular}{p{3in}|p{3in}}
         Stochastics (Individual Based Model)  & Squid Pop'n and Overfishing \\ \hline
         Could either do something related to a disease going through a pop'n of humble squid OR could do some kind of squid fishery thing ---- could do a disease spreading through a squid fishery with squeeze points and different disease dynamics -- something with changing the size of the space they can travel and also increasing the pop'n rapidly(ish) & Overfishing of other fish changes the squid carrying capacity - how does that change over time? \\
    \end{tabular}
    \caption{shoutout projects}
    % \label{tab:placeholder}
\end{table}

\textbf{Stochastics (Individual Based Model):}
\begin{table}[H]
    \centering
    \begin{tabular}{p{1.5in}|p{1.5in}|p{1.5in}|p{1.5in}}
         Resources  & Tasks & Attractiveness & Flaws \\ \hline
         literature, available code & modeling, code & link to stochastic modeling, accessible, would be cool to model swimming/the fluid motion of squid in a pop'n, cooler pictures & will require a concrete focus, lots of debugging \\
    \end{tabular}
    \caption{shoutout squid}
\end{table}


Squid Pop'n and Overfishing:
\begin{table}[H]
    \centering
    \begin{tabular}{p{1.5in}|p{1.5in}|p{1.5in}|p{1.5in}}
         Resources  & Tasks & Attractiveness & Flaws \\ \hline
         well studied i think, lots of resources  & modeling, code, analysis of model & Accessible and ODE-based (yay ODEs), much more already known analytical theory and graphing directions & will need some work to be novel, could get super complicated super quickly, mostly likely will spiral into a PDE model \\
    \end{tabular}
    \caption{shoutout fish}
    % \label{tab:placeholder}
\end{table}

potential ideas: population vs. temperature, population vs. populations of other fish?
\\$\Rightarrow$ I think we can combine these into one model! Would be super interesting (-Madi)
%also i fixed the tables #ImLearning
% i didnt know you could do arrows in latex

% changed the table bc it got long -sarah


\section{POTENTIAL RESOURCES} %tysm !!!!
%\begin{table}[H]
\begin{longtable}{p{2.0in}|p{1.0in}|p{2.5in}}
    \centering
    %\begin{tabular}{p{2.0in}|p{1.0in}|p{2.5in}}
         Title  & Author & Notes \\ \hline
         
        \href{https://onlinelibrary.wiley.com/doi/abs/10.1111/j.1365-2486.2010.02223.x}{Modelling climate-change-induced nonlinear thresholds in cephalopod population dynamics} & Andre et al., 2010 & maybe not relevant since we chose to do IBM rather than climate change population modelling \\ \\
        
        \href{https://link.springer.com/article/10.1007/s11160-007-9077-3}{The potential impacts of climate change on inshore squid: biology, ecology and fisheries} & Pecl \& Jackson, 2007 & also maybe not relevant since we chose to do IBM, but the paper says that maybe climate change will be bad for squids and have large effects (you might think it would be good bc they were repopulate faster, but they might need more food/oxygen so it'll affect everything else) \\ \\
        
        \href{https://www.annualreviews.org/content/journals/10.1146/annurev.ecolsys.36.102003.152644}{Individual-Based Modeling of Ecological and Evolutionary Processes} & Donald L. DeAngelis and Wolf M. Mooij, 2005 &  Review from `05 of individual-based models in ecology \\ \\

        \href{https://cdnsciencepub.com/doi/abs/10.1139/f93-015}{Growth Rate Variation and Larval Survival: Inferences from an Individual-Based Size-Dependent Predation Model} & Rice et al., 1993 & Use of a individual-based model to predict recruitment success of a cohort of fish based on length and growth rate\\ \\

        \href{https://www.degruyterbrill.com/document/doi/10.1515/9781400850624/html}{Individual-based Modeling and Ecology} & Volker Grimm and Steven F. Railsback, 2003 (book) & found a pdf online. book seems to be written for ecologists---mathematical detail seems a little sparse\\ \\ 
        
        \href{https://academic.oup.com/beheco/article/16/1/178/206370}{Density distribution and size sorting in fish schools: an individual-based model} & Charlotte K. Hemelrijk , Hanspeter Kunz, 2004 & IBM to see what spatial patterns exist in fish schools \\ \\ 

        \href{https://onlinelibrary.wiley.com/doi/abs/10.1046/j.1365-2419.1997.00037.x}{
An individual based model of fish school reactions: predicting antipredator behaviour as observed in nature} &RUNE VABØ, LEIF NØTTESTAD, 2003 & 2D model of herring to see internal school dynamics - might be useful to us since it talks about how the schools behave when attacked by a predator (a squid in our case) \\ \\

        \href{https://zslpublications.onlinelibrary.wiley.com/doi/abs/10.1111/j.1469-7998.1974.tb04385.x}{Experiments on whether schooling by their prey affects the hunting behaviour of cephalopods and fish predators} & Neill \& Cullen, 1974 & lab experiment to see how squid, cuttlefish, pike, and perch hunt when the school size is 1 / 6/ 20 \\ \\

        \href{https://www.frontiersin.org/journals/physiology/articles/10.3389/fphys.2017.00598/full}{Cephalopods as Predators: A Short Journey among Behavioral Flexibilities, Adaptions, and Feeding Habits} & Villanueva et al., 2017 & qualitative summary of how cephalopods hunt (i think) \\ \\

        \href{https://pmc.ncbi.nlm.nih.gov/articles/PMC2124471/}{Observations of wild hunting behaviour and bioluminescence of a large deep-sea, eight-armed squid, Taningia danae} & Kubodera et al., 2007 & how deep-sea squid use bioluminescence when hunting :3 \\ \\

        \href{https://link.springer.com/article/10.1007/s10164-020-00652-4}{Predation behaviour of the European squid Loligo vulgaris} & Castilla et al., 2020 & how does size, pattern, type of food, etc. affect how squid hunt and consume their prey? \\ \\

        \href{https://cdnsciencepub.com/doi/abs/10.1139/z85-003?casa_token=bqmjjmJ0cKEAAAAA:WwqN4ggP8UB-tploMwL8HRw3yJYGMCQLuByca4VsFl7IuEsnIVroR8C_c-qpgDbPYlPSr6QztWM}{Predatory behaviour of squid (Illex illecebrosus) feeding on surface swarms of euphausiids} & Nicol \& O'dor, 1985 & how canadian squid eat krill (apparently different from other observations?) edit: just saying they eat tentacle first, i thought this was common knowledge (maybe it changed since 1985) \\ \\

        \href{https://github-pages.ucl.ac.uk/rsd-engineeringcourse/ch02data/084Boids.html}{Bird-oid model code} & its a website & goes over how the Boid model works\\\\
    
    %\end{tabular} 
    \caption{shoutout literature}
    % \label{tab:placeholder}
%\end{table}
\end{longtable}


\section{TASKS (\& who's doing them)}
\begin{itemize}
    \item madi: build model
    \item cole:
    \begin{itemize}
        \item set up github / version control
        \item data analysis / statistics??
        \item making powerpoint fun
    \end{itemize}
    \item sarah: 
        \begin{itemize}
            \item writing?
            \item part of data analysis
            \item part of presentation
            \item is our model realistic
            \item background research
        \end{itemize}
        
\end{itemize}


\section{DELIVERABLES}

\subsection{Must-have's}
\begin{itemize}
    \item who did what
    \item a "diary" of experiential activities / making it better
    \begin{itemize}
        \item what did you try (2-3 sentences)
        \item how did you decide on it (be honest lol)
        \item talk about the before and after (how you think it improved)
    \end{itemize}
\end{itemize}

\subsection{functional code}
Deliverable details:
    \begin{itemize}
        \item the fish will be able to school
        \item the squid will hunt the fish
        \item fish will be hunted based on a) closeness to squid and b) distance from other fish
        \item boundary conditions (tank-like instead of periodic)
    \end{itemize}

HOW DO I FIX THE INDENT :( \\
Improving / sharpening:
\begin{itemize}
    \item "don't settle for broken code that sort of works"
    \item interrogate your code
    \item benchmarks
    \item use the benchmarks for timing (basically clean it up)
    \item ask the question, "what is a significant computation"
\end{itemize}

\subsection{talk/presentation ($\sim$15 mins), nobody does it like cole, DATE: NOV 26 2026}
Deliverable details:
\begin{itemize}
    \item 10-12 mins and 3-5 mins for Qs
    \item change your life or money back guaranteed
    \item will cover basics of Boid model
    \item will cover some basic statistics about the evolution of fish pop'n with squid hunting
    \item will cover fish schooling in regard to squid hunting
    \item playing with parameters and stuff
    \item cover how accurate it is to observed squid hunting
\end{itemize}

HOW DO I FIX THE INDENT \\ %wym ;-;
Improving / sharpening:
\begin{itemize}
    \item why does it matter?
    \item what does it say
    \item who is it speaking to
    \item ``there is a real theory of design out there"
    \item be aware of the wall of apathy
\end{itemize}

\subsection{other notes}
\begin{itemize}
    \item literature
    \begin{itemize}
        \item acknowledge there is a body of lit about the topic
        \item lead people to important parts
    \end{itemize}
\end{itemize}

\section{CORE MODEL}
\begin{itemize}
    \item FISH
    \begin{itemize}
        \item fish population / concentration
        \item how do they move (water current, brownian motion, etc.) (brownian :) done 11.12)
        \item amount of schooling / fish schooling mechanic (done 11.12)
    \end{itemize}
    \item SQUID
    \begin{itemize}
        \item jet propulsion? (madi doesn't know how to do this at all)
        \item choose fish based on multiple criteria? need mechanic
        \begin{itemize}
            \item how close to fish (done 11.12)
            \item how close fish is to other fish (mostly done (11.12)
        \end{itemize}
    \end{itemize}
\end{itemize}


\subsection{GOALS}
\begin{itemize}
    \item how much does schooling matter
    \begin{itemize}
        \item could investigate this in two respects: 1) vary the schooling/coupling factor for some subset (e.g., 10 or 20\%) of the fish and see if the fish that school more are more likely to survive; 2) vary the coupling factor for all fish and see if it increases overall survival rates. it'd be interesting to see if schooling behavior is important for overall survival rates, or only individual survival rates. my guess (hypothesis---i suppose) is that it'd only be important for the latter, not so much the former (the squid has to eat something, at the end of the day).
    \end{itemize}
    \item how much does the wait time between eating matter
    \item how much does squid-fish coupling matter?
    \item fish population / concentration over time
    \item realism??
\end{itemize}

\subsection{FUTURE WORK/SUPPLEMENTAL ACTIVITIES}
\begin{itemize}
    \item multiple squid / shoal of squid?
    \item water current / fish propulsion
\end{itemize}



\end{document}
